\ifdefined\gkver\else\def\gkver{student}\fi
\documentclass[\gkver]{gaokaozhenti}
\begin{document}

\chapter{2021年天津}

\begin{enumerate}
\item
%% number: 1
%% paperName: 2021年天津高考第 1 题 5 分
%% typeId: 1
%% type: 单选题
%% chapter: 其他
%% point: 物理思想
%% method: 无
%% score: 5
%% degree: 950
%% duplicateId: 0
%% body:
科学研究方法对物理学的发展意义深远，实验法、归纳法、演绎法、类比法、理想实验法等对揭示物理现象的本质十分重要。下列哪个成果是运用理想实验法得到的 \xzanswer{D}
\fourchoices[answer=D]
{牛顿发现“万有引力定律”}
{库仑发现“库仑定律”}
{法拉第发现“电磁感应现象”}
{伽利略发现“力不是维持物体运动的原因”}

%% answer: D
%% memo:
\memoanswer{
牛顿发现“万有引力定律”；库仑发现“库仑定律”；法拉第发现“电磁感应现象”，这些都是建立在大量的实验的基础上直接得出的结论；而伽利略发现“力不是维持物体运动的原因”，是在实验的基础上经过抽象推理得出的结论，即运用了理想实验法。

故选 D。
}

\item
%% number: 2
%% paperName: 2021年天津高考第 2 题 5 分
%% typeId: 1
%% type: 单选题
%% chapter: 光学
%% point: 光子说
%% method: 无
%% score: 5
%% degree: 850
%% duplicateId: 0
%% body:
光刻机是制造芯片的核心装备，利用光源发出的紫外线，将精细图投影在硅片上，再经技术处理制成芯片。为提高光刻机清晰投影最小图像的能力，在透镜组和硅片之间充有液体。紫外线进入液体后与其在真空中相比 \xzanswer{A}
\fourchoices[answer=A]
{波长变短}
{光子能量增加}
{频率降低}
{传播速度增大}

%% answer: A
%% memo:
\memoanswer{
紫外线进入液体后与真空相比，频率不变，传播速度减小，根据 $\lambda = \frac{v}{f}$ 可知波长变短；根据 $E = h\nu$ 可知，光子能量不变。

故选 A。
}

\item
%% number: 3
%% paperName: 2021年天津高考第 3 题 5 分
%% typeId: 1
%% type: 单选题
%% chapter: 交流电
%% point: 交流电
%% method: 无
%% score: 5
%% degree: 850
%% duplicateId: 0
%% body:
如图所示，闭合开关后，$R = 5 \UO$ 电阻两端的交流电压为 $u = 50\sqrt{2}\sin 10\pi t \UV$，电压表和电流表均为理想交流电表，则 \xzanswer{C}
\onepicture[w=4.2cm]{figs/03.png}
\fourchoices[answer=C]
{该交流电周期为 $0.02 \Us$}
{电压表的读数为 $100 \UV$}
{电流表的读数为 $10 \UA$}
{电阻的电功率为 $1 \UkW$}

%% answer: C
%% memo:
\memoanswer{
A．该交流电的周期

$T=\frac{2\pi}{\omega}=\frac{2\pi}{10\pi}\Us=0.2\Us$

B．电压表的读数为交流电的有效值，即

$U=\frac{50\sqrt{2}}{\sqrt{2}}\UV=50\UV$，选项 B 错误；

C．电流表的读数为

$I=\frac{U}{R}=\frac{50}{5}\UA=10\UA$，选项 C 正确；

D．电阻的电功率为

$P = IU = 10 \times 50\UW = 500\UW$，选项 D 错误。

故选 C。
}

\item
%% number: 4
%% paperName: 2021年天津高考第 4 题 5 分
%% typeId: 1
%% type: 单选题
%% chapter: 机械波
%% point: 波的图象
%% method: 无
%% score: 5
%% degree: 850
%% duplicateId: 0
%% body:
一列沿 $x$ 轴正方向传播的简谐横波，传播速度 $v = 10 \Ums$，$t = 0$ 时位于坐标原点的质点从平衡位置沿 $y$ 轴正方向运动，下列图形中哪个是 $t = 0.6 \Us$ 时的波形 \xzanswer{B}
\fourchoices[answer=B, ispicture=true, h=2.7cm]
{figs/04a.png}
{figs/04b.png}
{figs/04c.png}
{figs/04d.png}

%% answer: B
%% memo:
\memoanswer{
由图中可以看出该波的波长为 $\lambda = 4 \Um$，根据 $v = \frac{\lambda}{T}$ 可知该列波的周期为 $T = 0.4 \Us$。又因为 $t = 0$ 时位于坐标原点的质点从平衡位置沿 $y$ 轴正方向运动，当 $t = 0.6 \Us$ 时经历了 $1.5T$，所以此时位于坐标原点的质点从平衡位置沿 $y$ 轴负方向运动，结合图像可知 B 选项正确。

故选 B。
}

\item
%% number: 5
%% paperName: 2021年天津高考第 5 题 5 分
%% typeId: 1
%% type: 单选题
%% chapter: 曲线运动
%% point: 人造地球卫星
%% method: 无
%% score: 5
%% degree: 550
%% duplicateId: 0
%% body:
2021 年 5 月 15 日，天问一号探测器着陆火星取得成功，迈出了我国星际探测征程的重要一步，在火属上首次留下国人的印迹。天问一号探测器成功发射后，顺利被火星捕获，成为我国第一颗人造火星卫星。经过轨道调整，探测器先沿椭圆轨道 Ⅰ 运行，之后进入称为火星停泊轨道的椭圆轨道 Ⅱ 运行，如图所示，两轨道相切于近火点 P，则天问一号探测器 \xzanswer{D}
\onepicture[w=5.0cm]{figs/05.png}
\fourchoices[answer=D]
{在轨道 Ⅱ 上处于受力平衡状态}
{在轨道 Ⅰ 运行周期比在 Ⅱ 时短}
{从轨道 Ⅰ 进入 Ⅱ 在 P 处要加速}
{沿轨道 Ⅰ 向 P 飞近时速度增大}

%% answer: D
%% memo:
\memoanswer{
A．天问一号探测器在轨道 Ⅱ 上做变速圆周运动，受力不平衡，故 A 错误；

B．根据开普勒第三定律可知，轨道 Ⅰ 的半径大于轨道 Ⅱ 的半长轴，故在轨道 Ⅰ 运行周期比在 Ⅱ 时长，故 B 错误；

C．天问一号探测器从轨道 Ⅰ 进入 Ⅱ，做近心运动，需要的向心力要小于提供的向心力，故要在 P 点点火减速，故 C 错误；

D．在轨道 Ⅰ 向 P 飞近时，万有引力做正功，动能增大，故速度增大，故 D 正确。

故选 D。
}

\item
%% number: 6
%% paperName: 2021年天津高考第 6 题 5 分
%% typeId: 2
%% type: 多选题
%% chapter: 气体性质
%% point: 热力学第一定律
%% method: 无
%% score: 5
%% degree: 650
%% duplicateId: 0
%% body:
列车运行的平稳性与车厢的振动密切相关，车厢底部安装的空气弹簧可以有效减振，空气弹簧主要由活塞、气缸及内封的一定质量的气体构成。上下乘客及剧烈颠簸均能引起车厢振动，上下乘客时气缸内气体的体积变化缓慢，气体与外界有充分的热交换；剧烈颠簸时气缸内气体的体积变化较快，气体与外界来不及热交换。若气缸内气体视为理想气体，在气体压缩的过程中 \xzanswer{AC}
\fourchoices[answer=AC]
{上下乘客时，气体的内能不变}
{上下乘客时，气体从外界吸热}
{剧烈颠簸时，外界对气体做功}
{剧烈颠簸时，气体的温度不变}

%% answer: AC
%% memo:
\memoanswer{
AB．上下乘客时气缸内气体与外界有充分的热交换，即发生等温变化，温度不变，故气体的内能不变，体积变化缓慢，没有做功，故没有热交换，故 A 正确，B 错误；

CD．剧烈颠簸时气缸内气体的体积变化较快，且气体与外界来不及热交换，气体经历绝热过程，外界对气体做功，温度升高，故 C 正确，D 错误。

故选 AC。
}

\item
%% number: 7
%% paperName: 2021年天津高考第 7 题 5 分
%% typeId: 2
%% type: 多选题
%% chapter: 运动和力
%% point: 动量定理
%% method: 无
%% score: 5
%% degree: 750
%% duplicateId: 0
%% body:
一冲九霄，问鼎苍穹。2021 年 4 月 29 日，长征五号 B 遥二运载火箭搭载空间站天和核心舱发射升空，标志着我国空间站建造进入全面实施阶段。下列关于火箭的描述正确的是 \xzanswer{AB}
\fourchoices[answer=AB]
{增加单位时间的燃气喷射量可以增大火箭的推力}
{增大燃气相对于火箭的喷射速度可以增大火箭的推力}
{当燃气喷出火箭喷口的速度相对于地面为零时火箭就不再加速}
{火箭发射时获得的推力来自于喷出的燃气与发射台之间的相互作用}

%% answer: AB
%% memo:
\memoanswer{
A．增加单位时间的燃气喷射量，即增加单位时间喷射气体的质量，根据 $F\Delta t = mv$，可知可以增大火箭的推力，故 A 正确；

B．当增大燃气相对于火箭的喷射速度时，根据 $F\Delta t = \Delta mv$，可知可以增大火箭的推力，故 B 正确；

C．当燃气喷出火箭喷口的速度相对于地面为零时，此时火箭有速度，所以相对于火箭的速度不为零，火箭仍然受推力作用，仍然要加速，故 C 错误；

D．燃气被喷出的瞬间，燃气对火箭的反作用力作用在火箭上，使火箭获得推力，故 D 错误。

故选 AB。
}

\item
%% number: 8
%% paperName: 2021年天津高考第 8 题 5 分
%% typeId: 2
%% type: 多选题
%% chapter: 电场
%% point: 带电粒子的轨迹类问题
%% method: 无
%% score: 5
%% degree: 450
%% duplicateId: 0
%% body:
两个位于纸面内的点电荷产生电场的等势面如图中实线所示，相邻等势面间的电势差相等。虚线 MPN 是一个电子在该电场中的运动轨迹，轨迹与某等势面相切于 P 点。下列说法正确的是 \xzanswer{CD}
\onepicture[w=5.0cm]{figs/08.png}
\fourchoices[answer=CD]
{两点电荷可能是异种点电荷}
{A 点的电场强度比 B 点的大}
{A 点的电势高于 B 点的电势}
{电子运动到 P 点时动能最小}

%% answer: CD
%% memo:
\memoanswer{
A．根据电荷间等势面的分布情况可知两点电荷是同种电荷，又根据电子在该电场中的运动轨迹可判断电子一直受到排斥的力，故可知两点电荷为同种负电荷；故 A 错误；

B．根据等势面的疏密程度可以判断 A 点的电场强度比 B 点的小，故 B 错误；

C．因为两点电荷是同种负电荷，电场线指向负电荷，故可知 A 点的电势高于 B 点的电势，故 C 正确；

D．根据电子的运动轨迹和电场线的方向可知由 M 到 P 电场力做负功，由 P 到 N 电场力做正功；由 M 到 P 动能减小，由 P 到 N 动能增加，故电子运动到 P 点时动能最小，故 D 正确。

故选 CD。
}

\item
%% number: 9
%% paperName: 2021年天津高考第 9 题 6 分
%% typeId: 4
%% type: 实验题
%% chapter: 电路
%% point: 测电动势与内阻
%% method: 无
%% score: 6
%% degree: 650
%% duplicateId: 0
%% body:
某实验小组利用手机的录像功能拍下小球在斜面上做匀加速直线运动的过程。为便于记录小球各个时刻在斜面上的位置，将录像中时间间隔为 $T$ 的连续 7 幅画面合成到同一张图中，示意如图。依次测得小球各相邻位置间的距离为 $x_{1}$、$x_{2}$、$x_{3}$、$x_{4}$、$x_{5}$、$x_{6}$。

\onepicture[w=8.0cm]{figs/09a.png}

\begin{enumerate}
	\item
	写出小球在位置 1 的速度表达式 \tkanswer[3.0cm]{$\frac{x_{1}+x_{2}}{2T}$}。
	\item
	要求充分利用测量数据，写出小球运动过程中的加速度表达式 \tkanswer[3.6cm]{$\frac{(x_{4}+x_{5}+x_{6})-(x_{1}+x_{2}+x_{3})}{9T^{2}}$}。
	\item
	在测量小球相邻位置间距时由于实验者读数产生的误差是 \tkanswer{偶然} 误差。（填“偶然”或“系统”）
\end{enumerate}

随着智能手机的广泛应用，充电宝成为手机及时充电的一种重要选择。充电宝可以视为与电池一样的直流电源。一充电宝的电动势约为 $5 \UV$，内阻很小，最大放电电流为 $2 \UA$，某实验小组测定它的电动势和内阻。他们剥开充电宝连接线的外绝缘层，里而有四根导线，红导线为充电宝的正极，黑导线为充电宝的负极，其余两根导线空置不用，另有滑动变阻器 $R$ 用于改变电路中的电流，定值电阻 $R_{0} = 3 \UO$，两只数字多用电表 M、N，两表均为理想电表，并与开关 S 连成如图所示电路。

\onepicture[align=right, w=5.5cm]{figs/09b.png}

\begin{enumerate}
	\item
	图中测量电流的电表是 \tkanswer{N}，测量电压的电表是 \tkanswer{M}。（均填写字母“M”或“N”）
	\item
	调节滑动变阻器，测得多组 $I$、$U$ 数据，记录如下表，其中只有一个数据记录有误，审视记录的数据，可以发现表中第 \tkanswer{4} 次的记录数据有误。（填测量次数的序号）

	\begin{center}
	\begin{tblr}{|c|c|c|c|c|c|c|c|}
	\hline
	次数 & 1 & 2 & 3 & 4 & 5 & 6 & 7 \\
	\hline
	电流 $I$/$\UA[a]$ & 0.299 & 0.477 & 0.684 & 0.877 & 1.065 & 1.281 & 1.516 \\
	\hline
	电压 $U$/$\UV[a]$ & 4.970 & 4.952 & 4.932 & 4.942 & 4.894 & 4.872 & 4.848 \\
	\hline
	\end{tblr}
	\end{center}
	\item
	电路中接入 $R_{0}$ 可以达到下列哪个效果。\tkanswer{C}（填选项前的字母）

	A．使测电流的电表读数变化明显

	B．为了更准确地测量充电宝内阻

	C．避免使充电宝的放电电流过大

	D．减小测量电压的电表分流作用
\end{enumerate}

%% answer:
\jdanswer{
\begin{enumerate}
	\item
	$\frac{x_{1}+x_{2}}{2T}$

	\item
	$\frac{(x_{4}+x_{5}+x_{6})-(x_{1}+x_{2}+x_{3})}{9T^{2}}$

	\item
	偶然

\end{enumerate}
\begin{enumerate}
	\item
	N，M

	\item
	4

	\item
	C

\end{enumerate}
}
%% memo:
\memoanswer{
（1）
①匀变速直线运动中，中间时刻速度等于平均速度，所以 1 位置的速度为

$v_{1} = \frac{x_{1} + x_{2}}{2T}$

②题中要求充分利用数据，利用逐差法求解加速度，则

$x_{4} - x_{1} = 3a_{1}T^{2}$

$x_{5} - x_{1} = 3a_{2}T^{2}$

$x_{6} - x_{3} = 3a_{3}T^{2}$

解得加速度为

$a = \frac{a_{1} + a_{2} + a_{3}}{3} = \frac{(x_{4} + x_{5} + x_{6}) - (x_{1} + x_{2} + x_{3})}{9T^{2}}$

③读数产生的误差是人为因素造成的，属于偶然误差。

（点睛：一般来说，偶然误差是主观因素引起的误差，比如读数错误、操作错误等；系统误差是客观因素造成的，比如测量工具本身的精确度不够，测量方法带来的不可避免的误差等。）

（2）①器件 N 串联在电路中，测量干路电流，所以 N 为电流表；器件 M 一端接在开关上，另一端接在充电宝的负极，闭合开关后相当于并联在电源两端，所以 M 为电压表，测量路端电压；

②题中表格的电流 $I$ 不断增大，根据闭合电路欧姆定律 $E = U + Ir$ 可知路端电压 $U$ 不断减小，电压表的示数不断减小，所以第 4 次的记录数据有误；

③因为充电宝的内阻很小，所以电阻 $R_{0}$ 串联在电路中，起保护电路的作用，避免使充电宝的放电电流过大损坏电路中的器件。

故选 C。
}

\item
%% number: 10
%% paperName: 2021年天津高考第 10 题 14 分
%% typeId: 6
%% type: 计算题
%% chapter: 曲线运动
%% point: 平抛运动
%% method: 无
%% score: 14
%% degree: 650
%% duplicateId: 0
%% body:
一玩具以初速度 $v_{0}$ 从水平地面竖直向上抛出，达到最高点时，用遥控器将玩具内压缩的轻弹簧弹开，该玩具沿水平方向分裂成质量之比为 $1:4$ 的两部分，此时它们的动能之和与玩具从地面抛出时的动能相等。弹簧弹开的时间极短，不计空气阻力。求
\begin{enumerate}
	\item
	玩具上升到最大高度 $\frac{3}{4}$ 时的速度大小；
	\item
	两部分落地时速度大小之比。
\end{enumerate}

%% answer:
\jdanswer{
\begin{enumerate}
	\item
	$v = \frac{1}{2}v_{0}$

	\item
	$2$

\end{enumerate}
}
%% memo:
\memoanswer{
（1）设玩具上升的最大高度为 $h$，玩具上升到高度 $\frac{3}{4}h$ 时的速度大小为 $v$，重力加速度大小为 $g$，以初速度方向为正，整个运动过程有

$0 - v_{0}^{2} = -2gh$

玩具上升到最大高度 $\frac{3}{4}$ 有

$v^{2} - v_{0}^{2} = -2g\left(\frac{3}{4}h\right)$

两式联立解得

$v = \frac{1}{2}v_{0}$

（2）设玩具分开时两部分的质量分别为 $m_{1}$、$m_{2}$，水平速度大小分别为 $v_{1}$、$v_{2}$。依题意，动能关系为

$\frac{1}{2}m_{1}v_{1}^{2} + \frac{1}{2}m_{2}v_{2}^{2} = \frac{1}{2}(m_{1} + m_{2})v_{0}^{2}$

玩具达到最高点时速度为零，两部分分开时速度方向相反，水平方向动量守恒，有

$m_{1}v_{1} - m_{2}v_{2} = 0$

分开后两部分做平抛运动，由运动学关系，两部分落回地面时，竖直方向分速度大小为 $v_{0}$，设两部分落地时的速度大小分别为 $v_{1}'$、$v_{2}'$，由速度合成公式，有

$v_{1}' = \sqrt{v_{0}^{2} + v_{1}^{2}}$，$v_{2}' = \sqrt{v_{0}^{2} + v_{2}^{2}}$

结合 $m_{1}:m_{2} = 1:4$，解得

$\frac{v_{1}'}{v_{2}'} = 2$
}

\item
%% number: 11
%% paperName: 2021年天津高考第 11 题 16 分
%% typeId: 6
%% type: 计算题
%% chapter: 电磁感应
%% point: 电磁感应中的变加速
%% method: 无
%% score: 16
%% degree: 650
%% duplicateId: 0
%% body:
如图所示，两根足够长的平行光滑金属导轨 MN、PQ 间距 $L = 1 \Um$，其电阻不计，两导轨及其构成的平面均与水平面成 $\theta = 30^\circ$ 角，N、Q 两端接有 $R = 1 \UO$ 的电阻。一金属棒 ab 垂直导轨放置，ab 两端与导轨始终有良好接触，已知 ab 的质量 $m = 0.2 \Ukg$，电阻 $r = 1 \UO$，整个装置处在垂直于导轨平面向上的匀强磁场中，磁感应强度大小 $B = 1 \UT$。ab 在平行于导轨向上的拉力作用下，以初速度 $v_{1} = 0.5 \Ums$ 沿导轨向上开始运动，可达到最大速度 $v = 2 \Ums$。运动过程中拉力的功率恒定不变，重力加速度 $g = 10 \Umsq$。
\begin{enumerate}
	\item
	求拉力的功率 $P$；
	\item
	ab 开始运动后，经 $t = 0.09 \Us$ 速度达到 $v_{2} = 1.5 \Ums$，此过程中 ab 克服安培力做功 $W = 0.06 \UJ$，求该过程中 ab 沿导轨的位移大小 $x$。
\end{enumerate}
\onepicture[align=right, w=4.5cm]{figs/11.png}

%% answer:
\jdanswer{
\begin{enumerate}
	\item
	$P = 4 \UW$

	\item
	$x = 0.1 \Um$

\end{enumerate}
}
%% memo:
\memoanswer{
（1）在 ab 运动过程中，由于拉力功率恒定，ab 做加速度逐渐减小的加速运动，速度达到最大时，加速度为零，设此时拉力的大小为 $F$，安培力大小为 $F_{A}$，有

$F - mg\sin\theta - F_{A} = 0$

设此时回路中的感应电动势为 $E$，由法拉第电磁感应定律，有

$E = BLv$

设回路中的感应电流为 $I$，由闭合电路欧姆定律，有

$I = \frac{E}{R + r}$

ab 受到的安培力

$F_{A} = ILB$

由功率表达式，有

$P = Fv$

联立上述各式，代入数据解得

$P = 4 \UW$

（2）ab 从速度 $v_{1}$ 到 $v_{2}$ 的过程中，由动能定理，有

$Pt - W - mgx\sin\theta = \frac{1}{2}mv_{2}^{2} - \frac{1}{2}mv_{1}^{2}$

代入数据解得

$x = 0.1 \Um$
}

\item
%% number: 12
%% paperName: 2021年天津高考第 12 题 18 分
%% typeId: 6
%% type: 计算题
%% chapter: 磁场
%% point: 洛伦兹力
%% method: 无
%% score: 18
%% degree: 450
%% duplicateId: 0
%% body:
霍尔元件是一种重要的磁传感器，可用在多种自动控制系统中。长方体半导体材料厚为 $a$、宽为 $b$、长为 $c$，以长方体三边为坐标轴建立坐标系 $xyz$，如图所示。半导体中有电荷量均为 $e$ 的自由电子与空穴两种载流子，空穴可看作带正电荷的自由移动粒子，单位体积内自由电子和空穴的数目分别为 $n$ 和 $p$。当半导体材料通有沿 $+x$ 方向的恒定电流后，某时刻在半导体所在空间加一匀强磁场，磁感应强度的大小为 $B$，沿 $+y$ 方向，于是在 $z$ 方向上很快建立稳定电场，称其为霍尔电场，已知电场强度大小为 $E$，沿 $-z$ 方向。
\begin{enumerate}
	\item
	判断刚加磁场瞬间自由电子受到的洛伦兹力方向；
	\item
	若自由电子定向移动在沿 $+x$ 方向上形成的电流为 $I_{n}$，求单个自由电子由于定向移动在 $z$ 方向上受到洛伦兹力和霍尔电场力的合力大小 $F_{nz}$；
	\item
	霍尔电场建立后，自由电子与空穴在 $z$ 方向定向移动的速率分别为 $v_{nz}$、$v_{pz}$，求 $\Delta t$ 时间内运动到半导体 $z$ 方向的上表面的自由电子数与空穴数，并说明两种载流子在 $z$ 方向上形成的电流应满足的条件。
\end{enumerate}
\onepicture[align=right, w=7.5cm]{figs/12.png}

%% answer:
\jdanswer{
\begin{enumerate}
	\item
	自由电子受到的洛伦兹力沿 $+z$ 方向；

	\item
	$F_{nz} = e\left(\frac{I_{n}B}{neab} + E\right)$

	\item
	设 $\Delta t$ 时间内在 $z$ 方向上运动到半导体上表面的自由电子数为 $N_{n}$、空穴数为 $N_{p}$，则 $N_{n} = nacv_{nz}\Delta t$，$N_{p} = pacv_{pz}\Delta t$；霍尔电场建立后，半导体 $z$ 方向的上表面的电荷量就不再发生变化，则应 $N_{n} = N_{p}$，即在任何相等时间内运动到上表面的自由电子数与空穴数相等，这样两种载流子在 $z$ 方向形成的电流应大小相等、方向相反。

\end{enumerate}
}
%% memo:
\memoanswer{
（1）自由电子受到的洛伦兹力沿 $+z$ 方向；

（2）设 $t$ 时间内流过半导体垂直于 $x$ 轴某一横截面自由电子的电荷量为 $q$，由电流定义式，有

$I_{n} = \frac{q}{t}$

设自由电子在 $x$ 方向上定向移动速率为 $v_{nx}$，可导出自由电子的电流微观表达式为

$I_{n} = neabv_{nx}$

单个自由电子所受洛伦兹力大小为

$F_{\text{洛}} = ev_{nx}B$

霍尔电场力大小为

$F_{\text{电}} = eE$

自由电子在 $z$ 方向上受到的洛伦兹力和霍尔电场力方向相同，联立得其合力大小为

$F_{nz} = e\left(\frac{I_{n}B}{neab} + E\right)$

（3）设 $\Delta t$ 时间内在 $z$ 方向上运动到半导体上表面的自由电子数为 $N_{n}$、空穴数为 $N_{p}$，则

$N_{n} = nacv_{nz}\Delta t$

$N_{p} = pacv_{pz}\Delta t$

霍尔电场建立后，半导体 $z$ 方向的上表面的电荷量就不再发生变化，则应 $N_{n} = N_{p}$

即在任何相等时间内运动到上表面的自由电子数与空穴数相等，这样两种载流子在 $z$ 方向形成的电流应大小相等、方向相反。
}

\end{enumerate}

\end{document}
